%% title: Cochrane-Orcutt, Stochastic Regressors and Consistency
\documentclass[11pt]{article}
\usepackage{metricsstyle}

\begin{document}

\lessonheader{9}{Feasible GLS by Cochrane--Orcutt, stochastic regressors, and consistency}{}

\section{AR(1) errors with an unknown $\rho$}

The model of Lesson 8:
\[ y_{t} = x_{t}'\beta + u_{t}, \qquad u_{t} = \rho u_{t-1} + \varepsilon_{t}, \qquad
   \varepsilon_{t} \sim N(0, \sigma_{\varepsilon}^{2}), \qquad
   \E(\varepsilon_{t}\varepsilon_{s}) = 0 \ \ (t \neq s). \]
Lesson 8 assumed $\rho$ known. Now $\rho$ is unknown $\Rightarrow$ it has to be estimated.

\topic{Quasi-differencing again}
Write the model at $t$, and the model at $t-1$ multiplied by $\rho$:
\begin{align*}
  y_{t}        &= x_{t}'\beta + u_{t},\\
  \rho y_{t-1} &= \rho x_{t-1}'\beta + \rho u_{t-1}.
\end{align*}
Subtracting,
\[ y_{t} - \rho y_{t-1} = x_{t}'\beta - \rho x_{t-1}'\beta
   + \underbrace{u_{t} - \rho u_{t-1}}_{\varepsilon_{t}}, \]
\[ y_{t} = \rho y_{t-1} + x_{t}'\beta - \underbrace{\rho x_{t-1}'\beta}_{\rho\cdot\beta} + \varepsilon_{t}. \]
The product $\rho\cdot\beta$ introduces non-linearity: the model is no longer linear
in the parameters, so it cannot be estimated by one OLS regression.

\topic{Rewrite it in two ways}
\begin{align*}
  y_{t} - \rho y_{t-1} &= (x_{t} - \rho x_{t-1})'\beta + \varepsilon_{t} \tag{A}\\
  \text{or}\qquad y_{t} - x_{t}'\beta &= (y_{t-1} - x_{t-1}'\beta)\,\rho + \varepsilon_{t} \tag{B}
\end{align*}
Given $\rho$, (A) is linear in $\beta$; given $\beta$, (B) is linear in $\rho$
(its left side is $u_{t}$ and the bracket is $u_{t-1}$).

\topic{The Cochrane--Orcutt (CO) procedure}
Superscripts $(j)$ count the iterations.
\begin{flatnum}
  \item Start with $\rho^{(1)} = 0$ in (A): regress $y_{t}$ on $x_{t}$ (plain OLS)
        $\Rightarrow \beta^{(1)}$.
  \item Use $\beta^{(1)}$ in (B): with the residuals
        $e_{t}^{(1)} = y_{t} - x_{t}'\beta^{(1)}$, regress $e_{t}^{(1)}$ on
        $e_{t-1}^{(1)}$ $\Rightarrow \rho^{(2)}$.
  \item Go back to (A) with $\rho^{(2)}$: regress $y_{t} - \rho^{(2)}y_{t-1}$ on
        $x_{t} - \rho^{(2)}x_{t-1}$ $\Rightarrow \beta^{(2)}$.
  \item Use $\beta^{(2)}$ for $\beta$ in (B): regress $e_{t}^{(2)}$ on $e_{t-1}^{(2)}$
        $\Rightarrow \rho^{(3)}$.
  \item Use $\rho^{(3)}$ in (A), etc.
\end{flatnum}
until we find
\[ \bigl|\rho^{(j)} - \rho^{(j-1)}\bigr| < \delta, \qquad
   \bigl|\beta^{(j)} - \beta^{(j-1)}\bigr| < \delta, \qquad \delta \text{ small enough}. \]


\section{Stochastic regressors}

\[ y = X\beta + u, \qquad X \text{ is no longer fixed}. \]
$\rank(X) = K$, $X$ stochastic, $u \sim\ ?\,(0, \sigma^{2}I)$: the distribution of $u$ is
left unspecified.

\topic{Is OLS still unbiased?}
\[ b = (X'X)^{-1}X'y = \beta + (X'X)^{-1}X'u, \]
\[ \E(b) = \beta + \E\bigl[(X'X)^{-1}X'u\bigr]. \]
Now $(X'X)^{-1}X'$ is random too, so the expectation cannot be taken through it.
Condition on $X$ first and then apply the law of iterated expectations:
\[ \E\bigl[(X'X)^{-1}X'u\bigr] = \E\bigl[(X'X)^{-1}X'\,\underbrace{\E(u\mid X)}_{0}\bigr] = 0
   \qquad\Longrightarrow\qquad \E(b) = \beta, \]
provided
\[ \E(u \mid X) = 0. \]

\topic{$\E(u\mid X) = 0$ is stronger than $\E(u) = 0$}
\[ \E(u) = \E\bigl[\E(u\mid X)\bigr]: \qquad
   \E(u\mid X) = 0 \ \Longrightarrow\ \E(u) = 0, \qquad
   \E(u) = 0 \ \not\Longrightarrow\ \E(u\mid X) = 0. \]

\topic{Two conditions on $X$ and $u$}
\begin{align*}
  \E(u\mid X) &= 0 && \text{$X$ and $u$ are ``$\perp$''}\\
  \underbrace{\E(X'u)}_{\text{orthogonality condition}} &= 0 && \text{no correlation}
\end{align*}

\runin{Explanation.} The three ways $u$ can be unrelated to $X$ are nested, from
strongest to weakest:
\begin{flatlist}
  \item \emph{Independence}, $u \perp X$: the whole distribution of $u$ is the same
        whatever $X$ is (with $\E(u) = 0$).
  \item \emph{Mean independence}, $\E(u\mid X) = 0$: only the \emph{mean} of $u$ does
        not move with $X$; its variance, say, still may. This is what the line above
        means by ``$\perp$''. It is the condition that makes $b$ unbiased.
  \item \emph{Orthogonality} (no correlation), $\E(X'u) = 0$, i.e.\
        $\E(x_{t}u_{t}) = 0$: $u$ has no \emph{linear} relation with the regressors.
\end{flatlist}
Each implies the next: $\E(X'u) = \E\bigl[X'\,\E(u\mid X)\bigr] = 0$ by iterated
expectations, but not conversely. For instance, $u = x^{2} - 1$ with
$x \sim N(0,1)$ has $\E(xu) = \E(x^{3}) - \E(x) = 0$, yet $\E(u\mid x) = x^{2} - 1 \neq 0$.
The weaker orthogonality condition is not enough for unbiasedness, but it is what
consistency needs (below), and it is the moment condition behind IV and GMM in
Lesson~2.


\section{Large-sample theory: consistency}

\topic{The sample mean}
Let $y_{1}, y_{2}, \dots, y_{n}$ be iid $(\mu, \sigma^{2})$ and
\[ \bar y_{n} = \frac{1}{n}\sum_{i=1}^{n} y_{i}, \qquad
   \E(\bar y_{n}) = \mu, \qquad \Var(\bar y_{n}) = \frac{\sigma^{2}}{n}. \]
So
\[ \bar y_{n} \sim \Bigl(\mu, \frac{\sigma^{2}}{n}\Bigr), \qquad
   \frac{\sigma^{2}}{n} \to 0 \ \text{ as } n \to \infty, \]
\[ \bar y_{n} - \mu \sim \Bigl(0, \frac{\sigma^{2}}{n}\Bigr), \qquad
   \sqrt{n}\,(\bar y_{n} - \mu) \sim (0, \sigma^{2}), \qquad
   \sqrt{n}\,\Bigl(\frac{\bar y_{n} - \mu}{\sigma}\Bigr) \sim N(0, 1), \]
the last exactly when the $y_{i}$ are normal and approximately otherwise (the central
limit theorem, below).

\topic{Convergence in probability}
\[ |\bar y_{n} - \mu| < \delta \quad\Longleftrightarrow\quad
   \bar y_{n} - \delta < \mu < \bar y_{n} + \delta. \]
\begin{center}
\begin{tikzpicture}[scale=0.9]
  \ecaxes{8.4}{4}
  \node[ecnode,anchor=south east] at (0,4) {$\bar y_{n}$};
  \node[ecnode,below] at (8.4,0) {$n$};
  % band mu +/- delta
  \draw[red,dashed,line width=0.6pt] (0,2.6) -- (7.6,2.6);
  \draw[red,dashed,line width=0.6pt] (0,1.6) -- (7.6,1.6);
  \draw[line width=0.4pt,gray] (0,2.1) -- (7.6,2.1);
  \node[ecnode,left] at (0,2.1) {$\mu$};
  \node[ecnode,left,font=\footnotesize] at (0,2.6) {$\mu+\delta$};
  \node[ecnode,left,font=\footnotesize] at (0,1.6) {$\mu-\delta$};
  % sample path of the running mean
  \draw[line width=0.7pt]
    (0.6,3.4) -- (1.0,2.8) -- (1.5,1.25) -- (2.5,1.8) -- (3.4,2.4) -- (4.1,3.1)
    -- (4.6,2.5) -- (5.2,1.85) -- (5.7,1.95) -- (6.3,2.15) -- (7.4,2.1);
  \foreach \p in {(0.6,3.4),(1.5,1.25),(4.1,3.1),(5.2,1.85),(7.4,2.1)} \fill \p circle (1.4pt);
  \node[ecnode,above right] at (0.6,3.4) {$\bar y_{50}$};
  \node[ecnode,below] at (1.5,1.25) {$\bar y_{100}$};
  \node[ecnode,above] at (4.1,3.1) {$\bar y_{1000}$};
  \node[ecnode,below] at (5.2,1.5) {$\bar y_{5000}$};
  \node[ecnode,right] at (7.6,2.1) {$\bar y_{10{,}000}$};
\end{tikzpicture}
\end{center}
As $n$ grows the running mean may still leave the band $\mu \pm \delta$, but less
and less often:
\[ \lim_{n\to\infty} P\bigl(|\bar y_{n} - \mu| < \delta\bigr) = 1
   \qquad \text{for every } \delta > 0, \]
\[ \operatorname*{plim} \bar y_{n} = \mu \quad \text{as } n \to \infty:
   \qquad \bar y_{n} \text{ is consistent.} \]

\topic{Rules for plim}
Unlike expectations, probability limits pass through non-linear (continuous)
functions:
\begin{align*}
  \operatorname*{plim}(\bar y_{n} \pm \bar x_{n})
    &= \operatorname*{plim}\bar y_{n} \pm \operatorname*{plim}\bar x_{n},
    & \E(y \pm x) &= \E(y) \pm \E(x),\\
  \operatorname*{plim}\Bigl(\frac{\bar y_{n}}{\bar x_{n}}\Bigr)
    &= \frac{\operatorname*{plim}\bar y_{n}}{\operatorname*{plim}\bar x_{n}},
    & \E\Bigl(\frac{y}{x}\Bigr) &\neq \frac{\E(y)}{\E(x)},\\
  \operatorname*{plim}\Bigl(\frac{1}{\bar y_{n}}\Bigr)
    &= \frac{1}{\operatorname*{plim}\bar y_{n}}
\end{align*}
(for the ratios, as long as the plim in the denominator is not zero).

\topic{Law of large numbers}
\[ \bar y_{n} = \frac{1}{n}\sum_{i=1}^{n} y_{i}, \qquad
   \operatorname*{plim}\bar y_{n} = \mu \quad \text{since } \frac{\sigma^{2}}{n} \to 0. \]
\begin{center}
\begin{tikzpicture}[scale=0.9]
  \draw[line width=0.6pt] (-3.4,0) -- (3.4,0);
  \draw[dashed,line width=0.5pt] (0,0) -- (0,3.3);
  \node[ecnode,below] at (0,0) {$\mu$};
  \foreach \s in {1.3,0.75,0.4}
    \draw[line width=0.7pt,domain=-3.3:3.3,samples=120]
      plot (\x,{0.85*exp(-(\x)^2/(2*\s*\s))/\s});
  \node[ecnode,right,font=\footnotesize,align=left] at (1.0,1.9) {larger $n$:\\ tighter around $\mu$};
\end{tikzpicture}
\end{center}
The distribution of $\bar y_{n}$ piles up on $\mu$ as $n$ grows: the law of large
numbers.

\topic{Central limit theorem}
\[ \frac{\sqrt{n}\,(\bar y_{n} - \mu)}{\sigma} \ \xrightarrow{\ d\ }\ N(0, 1). \]


\section{Unbiasedness versus consistency}

\topic{Case 1: unbiased and consistent}
$\bar y_{n}$\quad\checkmark: $\E(\bar y_{n}) = \mu$ and $\operatorname*{plim}\bar y_{n} = \mu$.

\topic{Case 2: biased but consistent}
\[ s^{2} = \frac{e'e}{n - K}, \qquad \E(s^{2}) = \sigma^{2}. \]
Now suppose we have
\[ \hat\sigma^{2} = \frac{e'e}{n} = \frac{n - K}{n}\,s^{2}, \qquad
   \E(\hat\sigma^{2}) = \frac{n - K}{n}\,\sigma^{2} \neq \sigma^{2}
   \quad\text{biased}, \]
\[ \operatorname*{plim}\hat\sigma^{2} = \sigma^{2} \quad \text{as } n \to \infty, \]
since $(n - K)/n \to 1$.

\topic{Case 3: unbiased but inconsistent}
\[ \hat y_{n} = \frac{1}{2}y_{1} + \frac{1}{2}\Bigl(\frac{y_{2} + y_{3} + \cdots + y_{n}}{n - 1}\Bigr), \]
\[ \E(\hat y_{n}) = \frac{1}{2}\mu + \frac{1}{2}\Bigl(\frac{(n-1)\mu}{n-1}\Bigr)
   = \frac{1}{2}\mu + \frac{1}{2}\mu = \mu, \]
\[ \Var(\hat y_{n}) = \frac{\sigma^{2}}{4}
   + \underbrace{\frac{1}{4}\,\frac{\sigma^{2}}{n-1}}_{\to\,0 \text{ as } n\to\infty}
   \ \longrightarrow\ \frac{\sigma^{2}}{4} \neq 0. \]
\begin{center}
\begin{tikzpicture}[scale=0.9]
  \draw[line width=0.6pt] (-3.4,0) -- (3.4,0);
  \node[ecnode,below] at (0,0) {$\mu$};
  \foreach \s/\h in {1.25/1.0,1.05/0.95,0.9/0.9}
    \draw[line width=0.7pt,domain=-3.3:3.3,samples=120]
      plot (\x,{2.3*\h*exp(-(\x)^2/(2*\s*\s))});
  \node[ecnode,right] at (3.0,1.8) {doesn't shrink};
  \draw[-{Stealth[length=4pt]},line width=0.5pt] (3.0,1.8) -- (1.4,1.5);
\end{tikzpicture}
\end{center}
The spread stays at least $\sigma^{2}/4$, so $\hat y_{n} \not\to \mu$ in probability:
then we miss the target, however large $n$ is.

\section*{Next}
\[ y = X\beta + u \]
(the page stops here).


\begin{transcriptnotes}
  \item The set opens with a divider page carrying only ``class 9''; it is not
        reproduced. Page 7 holds only $y = X\beta + u$, the start of the next topic,
        and page 8 is blank.
  \item As in Lessons 3--8, $6^{2}$ is $\sigma^{2}$ (and $6_{\varepsilon}^{2}$ is
        $\sigma_{\varepsilon}^{2}$); $B$ for $\beta$ is typeset $\beta$; $k$ is
        typeset $K$; the regressor row $X_{t}$ is typeset $x_{t}$ as in Lesson 8.
        Transposes written ${}^{T}$ and ${}'$ are all typeset ${}'$.
        \textbf{Notation:} the population mean is written $u$ on pages 4--6; typeset
        $\mu$ so that it does not clash with the error $u$.
  \item \textbf{Notation:} the CO iterations are written $\rho'$, $\beta'$, $e_{t}'$,
        $\rho^{2}$, $\beta^{2}$, $e_{t}^{2}$, $\rho^{3}$, \ldots, and once
        $\rho_{2}$ (in $y_{t} - \rho_{2}y_{t-1}$); these are iteration counters, not
        primes or powers, and are typeset $\rho^{(1)}, \beta^{(1)}, e_{t}^{(1)},
        \rho^{(2)}, \ldots$. The list numbering, the name Cochrane--Orcutt for
        ``co procedure'', the definition of $e_{t}^{(j)}$ and the remark on what (A)
        and (B) are linear in are \textbf{added}.
  \item Page 2: \textbf{Corrected:} the AR(1) line reads
        $u_{t-1} = \rho u_{t-1} + \varepsilon_{t}$; it is $u_{t} = \rho u_{t-1} +
        \varepsilon_{t}$. The model is written $y_{t} = x_{t}'\beta + u$; typeset
        $u_{t}$. ``$/\rho$'' beside the $t-1$ equation is read as ``multiply by
        $\rho$''. \textbf{Corrected:} in the subtracted line
        $\ldots - \rho x_{t-1}' + \ldots$ the $\beta$ is missing; it is
        $\rho x_{t-1}'\beta$. ``$\rho\cdot\beta$ introduce non-linearity'' is kept;
        the sentence after it is \textbf{added}.
  \item Page 3: ``use $\beta^{2}$ in $\beta$ (B)'' is read as ``use $\beta^{(2)}$
        for $\beta$ in (B)''. The second stopping rule is written
        $|\beta^{j} - \beta^{j-1}| < \delta$ without the left bar; it is closed.
        ``$u \sim\ ?\,(0, \sigma^{2}I)$'' is kept with its question mark. ``then we
        can after apply \ldots'' is expanded into the iterated-expectations step;
        the page has $\E\bigl((X'X)^{-1}X'\E(u\mid X)\bigr)$ with a $0$ under the
        inner expectation. The arrows $\not\Rightarrow$ / $\Leftarrow$ between
        $\E(u) = \E(\E(u\mid X))$ and $\E(u\mid X) = 0$ are typeset as the two
        implications they state. The page writes ``$x, u$ are $\perp$'' beside
        $\E(u\mid X) = 0$ with the note ``explain claude''; the \textbf{Explanation}
        paragraph, which reads that $\perp$ as mean independence and orders it
        between independence and no correlation, is \textbf{added} in answer to it.
  \item Page 4: the iid assumption is written ``$y_{1}, \ldots, y_{n} \sim (u,
        6^{2})$''; ``iid'' is \textbf{added}. \textbf{Corrected:} the page has
        $\sqrt{n}(\bar y_{n} - \mu) \sim (0, n)$; multiplying by $\sqrt{n}$ multiplies
        the variance $\sigma^{2}/n$ by $n$, giving $(0, \sigma^{2})$, as the next line
        $\sqrt{n}(\bar y_{n} - \mu)/\sigma \sim N(0,1)$ requires. The remark that
        this is exact only under normality is \textbf{added}. The figure is redrawn:
        a running mean $\bar y_{n}$ against $n$, a line at $\mu$ and a red dashed
        band $\mu \pm \delta$, with $\bar y_{50}$ and $\bar y_{1000}$ above the band,
        $\bar y_{100}$ below it, $\bar y_{5000}$ just inside and $\bar y_{10{,}000}$
        on $\mu$. ``for every $\delta > 0$'' is \textbf{added}.
  \item Page 5: the plim rules and the comparison with $\E$ are laid out side by
        side; the page writes $\E(y/x) \neq \E(y)/\E(x)$ with a crossed-out equals.
        The sentence on continuous functions and the condition on the denominator are
        \textbf{added}. ``law of larger \#'' is typeset ``law of large numbers''. The
        figure (three bell curves on a common centre $\mu$, narrower each time) is
        redrawn. The CLT is written $\sim N(0,1)$; typeset $\xrightarrow{d}$.
  \item Page 6: \textbf{Corrected:} the page has
        $\operatorname*{plim}\hat\sigma = \sigma^{2}$; it is
        $\operatorname*{plim}\hat\sigma^{2} = \sigma^{2}$. ``since $(n-K)/n \to 1$''
        is \textbf{added}. The variance of $\hat y_{n}$ is checked: $\Var(\tfrac12
        y_{1}) = \sigma^{2}/4$ and $\Var\bigl(\tfrac12\cdot\text{mean of } n-1\bigr)
        = \tfrac14\sigma^{2}/(n-1)$, as written; the limit $\sigma^{2}/4$ is
        \textbf{added}. ``$\hat y_{n} \overset{?}{\to} \mu$ then we miss then
        target'' is typeset as $\hat y_{n} \not\to \mu$ and ``we miss the target''.
        The figure (bell curves that do not narrow, ``doesn't shrink'') is redrawn.
  \item Pages of the original, for tracing back: page 1 the divider; page 2 the AR(1)
        model with unknown $\rho$, the non-linearity, (A) and (B), the first CO
        steps; page 3 the rest of CO and the stopping rule, stochastic $X$,
        $\E(b)$, $\E(u\mid X)$ vs $\E(X'u)$; page 4 the sample mean, convergence in
        probability and its figure; page 5 plim rules, the LLN figure, the CLT;
        page 6 the three cases of bias vs consistency; page 7 $y = X\beta + u$;
        page 8 blank.
\end{transcriptnotes}

\end{document}
